\section{Benchmark and comparison with other approaches}\label{sec-benchmark}

\begin{objectif}
Measure, on a synthetic Totem market built contributor by contributor, what the structural detector adds to naive
checks and to the starting ideas: B-spline plus Gaussian mixture, small neural networks.
\end{objectif}

\subsection{Design}
The synthetic market is built from the contributors up, so that each mechanism of \cref{sec-diagnostic} appears as it
does in Totem, not as a shape added to a curve. Details are in \cref{app-synthetic}.
\begin{itemize}
\item \textbf{Four categories} (\cref{tab-categories}), eleven maturities from one month to five years, six dividend
periods (the rest of 2026, then 2027 to 2031), valuation date 14~October 2026.
\item \textbf{Contributors} have consistent views (\cref{hyp-coherent}): dividends made of a persistent factor and a
noise specific to each dividend, whose dispersion grows with maturity, and a smooth repo curve. Option contributors
price calls and puts at five strikes with their own forward and a noisy volatility surface. Panels shrink with
maturity and are nested; option and dividend panels are drawn independently.
\item \textbf{The consensus} is the mean after a median/median-absolute-deviation filter; dispersions and contributor
counts are published as Totem would. The controller's pipeline then runs unchanged.
\item \textbf{The FO marks} are the true dividends with a \pct{1} noise, and the true repo with a level error of
$3$\,bp (one standard deviation).
\item \textbf{The scenarios} (\cref{tab-scenarios}) are simulated for each test name with the same draw of
contributors: differences between scenarios are paired, and the signal-to-noise ratio of each injection is known
exactly.
\end{itemize}

\begin{table}[ht]
\centering
\caption{The four categories of the synthetic market.}\label{tab-categories}
\input{generated/\lang/tab_categories}
\end{table}

\begin{table}[ht]
\centering
\caption{Benchmark scenarios. The first six are anomalies of the Totem curve; for the last two, a Totem anomaly alarm
is not the desired outcome (see the text).}\label{tab-scenarios}
\small
\begin{tabular}{@{}>{\raggedright\arraybackslash}p{3.4cm}>{\raggedright\arraybackslash}p{7.5cm}
>{\raggedright\arraybackslash}p{3.3cm}@{}}
\toprule
Scenario & Injection & Size\\
\midrule
Spot mismatch & the option contributors price off another spot & $\pm\pct{0.03}$ to $\pm\pct{0.15}$\\
Error on one period & one period (2027 to 2029) wrong in the dividend consensus & $\pm\num{0.15}$ to $\pm\num{0.5}$ for
a spot of $100$\\
Ex-date disagreement & \pct{35} to \pct{70} of the option panel put the dividend closest to an expiry on its other side
& one dividend\\
Stale profile & allocation profile shifted so that a dividend crosses an expiry & $3$ to $20$ days beyond\\
Forward outlier & consensus calls shifted at one expiry & $3$ to $6$ standard deviations of carry noise\\
Biased long sub-panel & the contributors left after the largest drop-off are biased on later dividends &
$\pm\pct{8}$ to $\pm\pct{20}$\\
\midrule
FO repo mismark & FO repo shifted in parallel, Totem intact & $\pm15$ to $\pm50$\,bp\\
Gross/net bias & every Totem period multiplied by $1+\tau$ & $|\tau|$ from \pct{8} to \pct{18}\\
\bottomrule
\end{tabular}
\end{table}

\paragraph{Detectors.} Five detectors are compared:
(1) \emph{naive gap}: the largest $|q^{\tot}_k-q^{\fo}_k|$;
(2) \emph{naive shape}: the largest distance of the Totem repo from the interpolation of its two neighbours;
(3) \emph{B-spline plus Gaussian mixture on the curve}: the six coefficients of a cubic B-spline fitted to the Totem
curve (knots in $\ln T$), clustered by a Gaussian mixture per category whose number of components is chosen by BIC; the
score is the negative log-likelihood;
(4) \emph{the same on the Totem--FO difference};
(5) \emph{the structural detector} of \cref{sec-detection}.

\paragraph{Protocol.} \vNTrain{} clean names per category calibrate each detector to \vFPR{} false alarms (a score
quantile for the first four; $\kappa$ and the level scale for the structural detector). Then \vNTestPerCat{} new names
per category are simulated under each of the \vNScen{} scenarios, that is, \vNTest{} test names.

\subsection{Results}
\begin{table}[ht]
\centering
\caption{Share of names flagged, each detector being calibrated to \vFPR{} false alarms on clean history. The ``Clean''
row is the false alarm rate realised on new names. For the last two rows before the mean, a Totem anomaly alarm is
not the desired outcome.}\label{tab-detection}
\input{generated/\lang/tab_detection}
\end{table}

\begin{figure}[ht]
\centering\includegraphics[width=\linewidth]{figures/\lang/power.pdf}
\caption{Detection rate as a function of the signal-to-noise ratio $\lambda$ of the injected anomaly, over the six
Totem anomalies. The grey curve is the theoretical power of a known typical shape at the typical threshold
(\cref{cor-power}).}
\label{fig-power}
\end{figure}

\begin{figure}[ht]
\centering\includegraphics[width=\linewidth]{figures/\lang/attribution.pdf}
\caption{Mechanism named first by the structural detector, per scenario. Level, step, spike and season are the
families of typical shapes; ``roughness'' means only the global test fired; ``no alarm'' that the name was not
flagged.}
\label{fig-attribution}
\end{figure}

\paragraph{Reading the results.}
\begin{itemize}
\item \textbf{On average, the best methods are level.} Over the six Totem anomalies, the structural detector flags
\vMeanStruct{} of the names, the mixture on the difference \vMeanGmmDiff, on the curve \vMeanGmmCurve, the naive shape
check \vMeanShape{} and the naive gap \vMeanNaive.
\item \textbf{Forward outliers} are where carry space and the noise model pay off most: \vFwdStruct{} detected, against
\vFwdNaive{} for the naive gap, which cannot see a three- to six-sigma outlier under the long-end noise, and
\vFwdShape{} for the naive shape check.
\item \textbf{Spot mismatches} (\vSpotStruct) and \textbf{stale profiles} (\vSeasonStruct) are easy for every method;
the structural detector adds the explanation (\cref{fig-attribution}).
\item \textbf{Errors on one period} (\vBucketStruct, against \vBucketNaive{} for the naive gap) are limited by their
size: $\num{0.15}$ to $\num{0.5}$ for a spot of $100$ is small against the dispersion of distant dividends, and
\cref{fig-power} shows that the misses have a low signal-to-noise ratio.
\item \textbf{Ex-date disagreements} are detected in \vTimingStruct{} of cases; the naive shape check (\vTimingShape) is
competitive, because these spikes are large in repo.
\item \textbf{The biased long sub-panel} escapes every method (\vPanelStruct{} for the structural detector,
\vPanelNaive{} for the naive gap): the bias builds up smoothly and blends into a repo slope. This is
\cref{thm-nonid}, not a tuning problem.
\item \textbf{Genuine FO mismarks} are not taken for bad Totem curves: \vFoStruct{} alarms (the clean rate), and
\vFoMismark{} are flagged as possible FO mismarks, against \vFoClean{} of clean names. The naive gap and the mixture on
the difference raise a bad-curve alarm on \vFoNaive{} and \vFoGmmDiff{} of them (\cref{prop-gmm-fo}). The mismark test
finds \vFoMismarkEuStock{} of the errors on European stocks and \vFoMismarkIndex{} on the index, but only
\vFoMismarkHtb{} on hard-to-borrow names: contributors' repo views there differ by about $60$\,bp, so an FO difference
of $15$ to $50$\,bp stays within the uncertainty of the consensus itself.
\item \textbf{The gross/net bias} is flagged in \vGrossStruct{} of cases, mostly as steps at the ex-dates of discrete
payers; for an index dividend stream it looks like a repo slope (\cref{prop-signatures}, (vi)).
\item \textbf{Power.} For a signal-to-noise ratio between $4$ and $6$, the structural detector finds \vPowStructE{} of
the anomalies, and above $6$ from \vPowStructF{} to \vPowStructH, in line with \cref{cor-power}. Below $3$, the naive
gap and the mixture on the difference detect more (\vPowGmmDiffC{} against \vPowStructC{} between $2$ and $3$): they
also react to the part of an anomaly that looks like a smooth repo difference, which the structural detector
deliberately leaves to the FO mismark test. The same property makes them flag genuine FO mismarks.
\item \textbf{Explanation.} When the structural detector flags a Totem anomaly, its first explanation names the
injected mechanism and its position in \vAttr{} of cases.
\end{itemize}
Results by category, areas under the ROC curve and realised false alarm rates are in \cref{app-results}.

\subsection{What becomes of the starting ideas}
\paragraph{B-spline plus Gaussian mixture.} The starting idea summarised each curve by its B-spline coefficients, by
region and category, and clustered these coefficients with a Gaussian mixture. Two properties limit this approach.

\begin{proposition}[What the coefficients do not see]\label{prop-leverage}
Let $Z\in\M_{N,n_s}(\R)$ be the matrix of the $n_s$ B-spline functions at the pillars and $H=Z(\tr{Z}Z)^{-1}\tr{Z}$ the
associated projection matrix. The least-squares fit of a unit spike $e_k$ has energy $\norm{He_k}^2=h_{kk}$, with
$\sum_kh_{kk}=n_s$. A fraction $1-h_{kk}$ of the spike's energy stays in the residual and escapes every feature computed
from the coefficients.
\end{proposition}
\begin{proof}
$H$ is symmetric and idempotent, so $\norm{He_k}^2=\tr{e_k}He_k=h_{kk}$, and $\sum_kh_{kk}=\operatorname{tr}H=n_s$. The
coefficients $(\tr{Z}Z)^{-1}\tr{Z}e_k$ determine $He_k$ and nothing else.
\end{proof}

For the benchmark's six cubic B-splines and eleven pillars:
\begin{center}\input{generated/\lang/tab_leverage}\end{center}
At interior pillars, the coefficients capture only \vLevMin{} to \vLevMax{} of a spike's energy, and the spline spreads
what it captures over several coefficients. A single-pillar anomaly, the most frequent pattern, is diluted before the
clustering even starts.

\begin{proposition}[Why a mixture on the difference flags FO errors]\label{prop-gmm-fo}
Take as features the coefficients $c=Z^+(q^{\tot}-q^{\fo})$, where $Z^+$ is the pseudo-inverse of $Z$. A parallel shift
$u$ of the FO repo moves them by $-u\,Z^+\one$, a fixed direction. If the mixture is fitted on clean history, where
Totem and the FO agree up to noise, its density decreases along this direction: its alarm rate on FO mismarks tends to
$1$ as $|u|$ grows, whatever the quality of the Totem curve.
\end{proposition}
\begin{proof}
$c$ is linear in $q^{\tot}-q^{\fo}$, which moves by $-u\one$. Every Gaussian component has finite variance along
$Z^+\one$: the mixture density, and with it the likelihood score, goes to zero along this direction.
\end{proof}
The structural detector does not have this defect, because $X\beta$ contains the level difference: it is removed from
every test and published separately. The benchmark confirms both statements (\vFoGmmDiff{} against \vFoStruct).

\paragraph{Small neural networks per anomaly type.}
\begin{proposition}[A network can only learn the matched filter]\label{prop-nn}
After whitening and projection, $\eh\sim\N(0, \Pi)$ under $H_0$ and $\eh\sim\N(a\gt, \Pi)$ under $H_1$, for a known shape
$g$ and size $a$. Then
\begin{equation}\label{eq-nn}
\ln\frac{\mathbb P(H_1\mid\eh)}{\mathbb P(H_0\mid\eh)}=a\sprod{\gt}{\eh}-\frac{a^2}{2}\norm{\gt}^2
+\ln\frac{\mathbb P(H_1)}{\mathbb P(H_0)},
\end{equation}
an affine function of the matched filter $\sprod{\gt}{\eh}$. A classifier trained by cross-entropy on enough data
converges to this posterior probability; with an unknown size and sign, one recovers the statistic $|z_g|$ of
\cref{thm-glrt}.
\end{proposition}
\begin{proof}
The two Gaussian distributions have the same covariance: the quadratic terms of their ratio cancel. Cross-entropy is
minimised by the true posterior probability; maximising \eqref{eq-nn} over $a$ gives $\sprod{\gt}{\eh}^2/(2\norm{\gt}^2)$.
\end{proof}

For the mechanisms one can name, a network can therefore at best learn the matched filter, which the structural
detector writes down directly, with an exact false alarm rate and a known power.

\paragraph{Where these tools remain useful.}
\begin{itemize}
\item \textbf{Penalised splines in the null model}, if $\beta_0+\beta_1\ln T$ turns out too rigid on real data.
\item \textbf{Gaussian mixtures on the detector's outputs} ($\hat\beta_0$, $\hat\beta_1$, the $z_g$, $\norm{\eh}^2$, the
dispersion ratios), to learn the usual regimes by region and category in a low-dimensional space where every
coordinate has a meaning.
\item \textbf{Region $\times$ category} to pool the calibration of $\kappa$ and of the level scale.
\item \textbf{Neural networks}, once there is a history of labelled incidents and patterns the typical shapes do not
cover; they should then be trained on the residuals $\eh$ and on synthetic injections, not on raw curves.
\end{itemize}

\subsection{Limits}
\begin{attention}
The synthetic market reproduces the mechanisms by construction; it does not show that real Totem data contain only
these. Real data will also include stale contributors, corporate actions and special dividends, American exercise in
single-stock quotes, and input errors. The scale $\kappa$ absorbs a misspecified noise level, not a misspecified
structure, and the injected sizes, though plausible, are not calibrated to any history. The benchmark shows that the
method works when the mechanisms are present and is calibrated when they are absent; validation on twelve to
twenty-four months of real history, with labelled incidents, is the next step.
\end{attention}

\begin{retenir}
\begin{itemize}[leftmargin=12pt]
\item At equal false alarm rates, the structural detector detects as much as the best alternative on average, and
clearly more on forward outliers.
\item It does not take an FO mismark for a bad curve, unlike a Gaussian mixture on the Totem--FO difference, and it
explains its alarms.
\item B-splines, mixtures and neural networks remain useful, but downstream of the structural model, not on raw curves.
\end{itemize}
\end{retenir}
